\documentclass{situationist-preprint}

\title{The Demandingness of Morality}
\msid{ETH-28-a}
\msfamily{ETH-28}
\msdiscipline{Normative and Applied Ethics}
\mstype{Principal manuscript}
\mslean{lean/Philosophy/ETH/ETH\_28.lean}
\mskeywords{demandingness; personal prerogative; beneficence; rescue; supererogation}

\begin{document}
\maketitle

\begin{abstract}
Morality is moderately demanding. We combine three results: the duty of beneficence requires affluent agents to give a tenth of their income (ETH-06); the duty to rescue requires rescues up to a cumulative cost of a tenth of income per year, beyond which rescue is supererogatory (ETH-18); and supererogatory acts, bounded below by decency and above by sainthood, are not required (ETH-14). We prove that the total requirement lies between a tenth and a fifth of the agent's income, that at least four fifths is left to a personal prerogative, and that there is room for supererogation above the line of decency. The proof is formalized in Lean.
\end{abstract}

\section{The question}

Objections to consequentialism from demandingness hold that a morality requiring agents to promote the good impartially would leave them no life of their own \citep{Scheffler1982}; \citet{Kagan1989} replied that the limits of ordinary morality cannot be defended. A dispute about how much morality demands is, however, not settled by the choice of framework, since every deontic theory has a consequentialist counterpart (ETH-02). It is settled by the content of the duties. We take the content from the three manuscripts of this cluster and compute their sum.

\section{Framework}

For an agent with annual income $I$, all costs are measured in income-equivalents for the year.

\begin{assumption}[Beneficence, ETH-06]\label{a:ben}
The duty of beneficence requires a tenth of income.
\end{assumption}

\begin{assumption}[Rescue, ETH-18]\label{a:res}
The rescues required in a year cost in total at most a tenth of income.
\end{assumption}

\begin{assumption}[Supererogation, ETH-14]\label{a:sup}
Supererogatory acts are not required: they contribute nothing to the requirement.
\end{assumption}

\begin{assumption}\label{a:saint}
The saintly bound on self-sacrifice does not exceed the agent's income.
\end{assumption}

\begin{definition}\label{d:required}
The \emph{total requirement} is the sum of the costs required by beneficence, rescue and supererogation. It is the \emph{line of decency}: an agent who meets it acts decently, and one who falls short does not.
\end{definition}

\begin{definition}\label{d:prerogative}
The \emph{personal prerogative} is the part of income not subject to requirement: income minus the total requirement.
\end{definition}

The two duties are additive because they are distinct. Beneficence is a share of a collective task; rescue is owed to particular persons in emergencies; neither draws on the other (ETH-06, ETH-18).

\section{The result}

\begin{theorem}\label{t:moderate}
Given Assumptions~\ref{a:ben}--\ref{a:sup}, the total requirement lies between a tenth and a fifth of income, and the personal prerogative is at least four fifths of income.
\end{theorem}

\begin{proof}
The requirement is $I/10$ for beneficence, at most $I/10$ for rescue and $0$ for supererogation, so it is at least $I/10$ and at most $I/5$. The prerogative is $I$ minus at most $I/5$, which is at least $4I/5$.
\end{proof}

\begin{corollary}\label{c:room}
Given also Assumption~\ref{a:saint}, if the saintly bound exceeds a fifth of income, there is a level of sacrifice above the line of decency and within the saintly bound.
\end{corollary}

\begin{proof}
The saintly bound itself is such a level, since by Theorem~\ref{t:moderate} the line of decency is at most a fifth of income.
\end{proof}

The theorem gives the sense in which morality is moderately demanding. It is demanding: an affluent agent who gives nothing and rescues no one falls below decency, and the requirement is never less than a tenth of income. It is moderate: the requirement never exceeds a fifth, whatever the sequence of emergencies or the level of others' compliance. Between the line of decency and the saintly bound lies the space of supererogation, which ETH-14 describes; beyond the requirement, what the agent does with four fifths of her income is hers to decide, including giving priority to those close to her (ETH-19).

\begin{objection}
The figures are those of affluent agents in ordinary circumstances. In a catastrophe, morality may demand more.
\end{objection}

\begin{reply}
The result is conditional on the assumptions, and a catastrophe that raised the cost of meeting basic need beyond the bound on which ETH-06 relies would change the beneficence share. What does not change is the structure: both duties are bounded fractions of income, so their sum is bounded, and a prerogative remains.
\end{reply}

\section{Formalization}

Theorem~\ref{t:moderate} and Corollary~\ref{c:room} are formalized in \nolinkurl{Philosophy/ETH/ETH_28.lean}. The three required costs and the saintly bound are primitive functions of income, and Assumptions~\ref{a:ben}--\ref{a:saint} are axioms; the results of ETH-06 and ETH-18 enter as premises here, not as imported theorems.

\begin{lstlisting}
theorem moderate (I : Nat) :
    I / 10 ≤ required I ∧ required I ≤ I / 5 ∧ 4 * I / 5 ≤ prerogative I
\end{lstlisting}

\begin{lstlisting}
'Philosophy.ETH_28.moderate' depends on axioms:
  [propext, Classical.choice, Quot.sound,
   Philosophy.ETH_28.beneficence, Philosophy.ETH_28.beneficence_tenth,
   Philosophy.ETH_28.rescue, Philosophy.ETH_28.rescue_tenth,
   Philosophy.ETH_28.supererogation, Philosophy.ETH_28.supererogation_zero]
\end{lstlisting}

\section{Conclusion}

The combined duties of beneficence, rescue and supererogation require at least a tenth and at most a fifth of an affluent agent's income, leaving at least four fifths to a personal prerogative and room for supererogation above decency. Morality is moderately demanding.

\begin{thebibliography}{9}\small
\bibitem[Kagan(1989)]{Kagan1989} Kagan, S. (1989). \emph{The Limits of Morality}. Clarendon Press.
\bibitem[Scheffler(1982)]{Scheffler1982} Scheffler, S. (1982). \emph{The Rejection of Consequentialism}. Clarendon Press.
\end{thebibliography}

\end{document}
